
Se comprobó que los Patrones Binarios Locales son un método efectivo para la 
extracción de características espaciales. Al combinar los histogramas LBP con un 
clasificador supervisado K-NN, se logró una exactitud máxima del $89.30\%$, demostrando 
que la frecuencia de estos micropatrones sirve como una ``firma´´ para clasificar 
texturas de diferentes materiales.

Al comparar los tres modelos evaluados, se observó un comportamiento no esperado: 
el rendimiento del modelo disminuyó conforme se aumentó el radio de búsqueda 
(obteniendo $89.30\%$ para $R=1$, $86.77\%$ para $R=2$, y $86.42\%$ para $R=3$). 
Esto indica que, para el conjunto de datos \texttt{KTH\_TIPS}, las características 
que mejor distinguen a las clases residen en las texturas más pequeñas, como los poros 
de la esponja o los granos de la lija, más que en las estructuras de mayor tamaño. 

La experimentación con $P=24$ dejó en evidencia el costo computacional exponencial 
del algoritmo LBP estándar, el cual generó vectores de $2^{24}$ dimensiones que 
desbordaron la memoria. Se concluye que el uso de métodos de reducción como LBP 
\textit{'uniform'} es indispensable en parámetros altos, ya que permitió reducir la 
dimensionalidad a solo 26 características manteniendo una exactitud de ($86.42\%$).

Se observó que los algoritmos de clasificación basados en árboles subyacentes en 
K-NN colapsan ante la alta dimensionalidad. La transición a una búsqueda de 
distancias por fuerza bruta y el ajuste del tipo de dato a punto flotante de 32 
bits fueron cruciales para resolver los cuellos de botella durante la fase de evaluación.

La configuración inicial de $P=8, R=1$ resultó ser la más óptima para este problema 
en específico, logrando una entre precisión de clasificación, bajo consumo de recursos 
computacionales y rapidez de ejecución algorítmica.

\QED